% ===== tr-02 块1：第 3 章 pp.5--7（原文页 5--7）=====
% [原文 p.5]
现在我们可以利用这一点来研究 $f$ 与 $\check f$ 的 Leray 谱序列。

我们再加一条假设：$X$ 与 $\check X$ 是单连通的。于是特别地 $B$ 是单连通的。再设 $B$ 是三维流形。这样我们就得到 $f$ 的 Leray 谱序列中的 $E_2$ 项：

\begin{center}
\begin{tikzpicture}[font=\small, x=3.2cm, y=1.35cm]
\node at (0,3) {$\R$};
\node at (1,3) {$0$};
\node at (2,3) {$0$};
\node at (3,3) {$\R$};
\node at (0,2) {$H^0(B,R^2f_*\R)$};
\node at (1,2) {$H^1(B,R^2f_*\R)$};
\node at (2,2) {$H^2(B,R^2f_*\R)$};
\node at (3,2) {$H^3(B,R^2f_*\R)$};
\node at (0,1) {$H^0(B,R^1f_*\R)$};
\node at (1,1) {$H^1(B,R^1f_*\R)$};
\node at (2,1) {$H^2(B,R^1f_*\R)$};
\node at (3,1) {$H^3(B,R^1f_*\R)$};
\node at (0,0) {$\R$};
\node at (1,0) {$0$};
\node at (2,0) {$0$};
\node at (3,0) {$\R$};
\end{tikzpicture}
\end{center}

% [Tutor Note: 谱序列图按 $E_2^{p,q}=H^p(B,R^qf_*\R)$ 排：横向为底方向次数 p=0..3，纵向为纤维方向次数 q=3（最上）到 0（最下）。最上一行与最下一行都是 R,0,0,R，因为 $R^0f_*\R=R^3f_*\R=\R$，而 B 单连通的三维流形满足 $b_1=b_2=0$，故只剩 $H^0$ 与 $H^3$ 两个 R。]

% [Tutor Note: 由 $H^1(X,\R)=0$ 推出 $H^0(B,R^1f_*\R)=0$ 的省略步骤：$B$ 单连通时 $E_2^{2,0}=H^2(B,\R)=0$，于是 $(0,1)$ 项的 $d_2$ 靶为零，$E_2^{0,1}$ 只能整体作为 $H^1(X,\R)$ 的一个滤过商出现，故为零；同理 $H^5(X,\R)=0$ 逼出 $E_2^{3,2}=H^3(B,R^2f_*\R)=0$（$p>3$ 的项全为零，它无路可去）。]

由于 $X$ 单连通，$H^1(X,\R)=H^5(X,\R)=0$，由此可断定 $H^0(B,R^1f_*\R)=H^3(B,R^2f_*\R)=0$。同样的论证对 $\check f$ 也成立，再由式 (3.1) 便得到 $H^0(B,R^2f_*\R)\cong H^0(B,R^1\check f_*\R)=0$，并类似地有 $H^3(B,R^1f_*\R)=0$。最后，考察谱序列中可能非零的映射：

\begin{center}
\begin{tikzpicture}[font=\small, x=3.2cm, y=1.45cm]
\node (r1) at (0,3) {$\R$};
\node at (1,3) {$0$};
\node at (2,3) {$0$};
\node at (3,3) {$\R$};
\node at (0,2) {$0$};
\node at (1,2) {$H^1(B,R^2f_*\R)$};
\node (h22) at (2,2) {$H^2(B,R^2f_*\R)$};
\node at (3,2) {$0$};
\node at (0,1) {$0$};
\node (h11) at (1,1) {$H^1(B,R^1f_*\R)$};
\node at (2,1) {$H^2(B,R^1f_*\R)$};
\node at (3,1) {$0$};
\node at (0,0) {$\R$};
\node at (1,0) {$0$};
\node at (2,0) {$0$};
\node (r4) at (3,0) {$\R$};
\draw[->] (r1.south east) -- node[midway, above] {$d_1$} (h22.north west);
\draw[->] (h11.south east) -- node[midway, above] {$d_2$} (r4.north west);
\end{tikzpicture}
\end{center}

% [Tutor Note: 这两条箭头的标注需按端点理解：按标准双度约定 $d_r:E_r^{p,q}\to E_r^{p+r,q-r+1}$，从 $(p,q)=(0,3)$ 指向 $(2,2)$、即从 $H^0(B,R^3f_*\R)$ 到 $H^2(B,R^2f_*\R)$ 的映射属于 $d_2$ 型；原文图中标作 $d_1$，译文照原文标注，读者以箭头起点/终点为准。第二条由 $H^1(B,R^1f_*\R)$ 到 $H^3(B,R)=\R$ 的箭头则是标准的 $d_2$。]

我们还需要一条进一步的假设，它在 $f_0$ 与 $\check f_0$ 的对偶关系下同样是自然的：$f$ 与 $\check f$ 都拥有截面。（其实，在 $\R$ 上工作时，我们只需 $X$ 与 $\check X$ 上存在上同调类，使其在 $f$ 或 $\check f$ 的某个纤维上的取值非零。）截面与每个纤维都非平凡地相交，因而给出 $R^3f_*\R$ 的一个截面。由于这样一个截面同时代表 $X$ 上的一个上同调类，$d_1$ 必是零映射。类似地，$f$ 的纤维不可能在同调意义下平凡，因为它与截面非平凡地相交，于是映射 $d_2$ 必为零。因此谱序列在 $E_2$ 处退化。

特别地，我们得到 $\R^{h^{1,1}}\cong H^2(X,\R)\cong H^1(B,R^1f_*\R)\cong H^1(B,R^2\check f_*\R)$ 以及 $\R^{h^{2,2}}\cong H^4(X,\R)\cong H^2(B,R^2f_*\R)\cong H^2(B,R^1\check f_*\R)$，其中 $h^{p,q}$ 是 $X$ 的 Hodge 数。于是 $\check X$ 的第三个 Betti 数是 $2+h^{1,1}+h^{2,2}=2(1+h^{1,1})$，由此可见 $h^{1,1}(X)=h^{1,2}(\check X)$、$h^{1,2}(X)=h^{1,1}(\check X)$。

% [Tutor Note: 记号 $\R^{h^{1,1}}$ 表示维数为 $h^{1,1}$ 的实向量空间；这里用 Betti 数配平：$\check X$ 的 $b_3=2+h^{1,1}+h^{2,2}=2(1+h^{1,1})$，而镜像对称要求 $h^{1,1}$ 与 $h^{1,2}$ 互换，$b_3=2+2h^{1,1}$ 在互换下不变，这正是三维情形 Hodge 数交换的自洽性检验。]

因此，除了一些当然最终必须被论证的假设之外，至少在三维情形下，Hodge 数为何会在对偶之下互换，是清楚的。

更一般地，在任意维数下，人们可以期望 $\R$-simple 性蕴含 $\dim_{\R}H^p(B,R^qf_*\R)=h^{p,q}$，于是 Hodge 数更一般的交换也就清楚了。相关结果见定理 3.8.10。

\begin{remark}{3.1.1}
该论证可以在 $\Z$ 上细化，从而对镜像对称下整系数上同调的行为作出新的预测。在 [16] 的定理 3.10 中（同样是在三维情形下）证明了：若 $f$ 与 $\check f$ 是 $\Z$-simple
% [原文 p.6]
且 $\Q/\Z$-simple 的，$f$ 与 $\check f$ 都有截面，并且 $H^1(X,\Z)=0$，则
\begin{align*}
H^{even}(X,\Z[1/2]) &\cong H^{odd}(\check X,\Z[1/2])\\
H^{odd}(X,\Z[1/2]) &\cong H^{even}(\check X,\Z[1/2]).
\end{align*}
\end{remark}

% [Tutor Note: $\Z[1/2]$ 表示把 2 可逆化之后的系数环，即只允许分母为 2 的幂；选择它正是为了绕开论证中的二挠（two-torsion）困难，故原文随后说「很可能上述同构在 Z 上成立」。]

论证中存在与二挠有关的问题，但上述同构很可能在 $\Z$ 上也成立。关于后一猜想的证据见 [4]。

不必再猜了。现在我们认真对待特殊拉格朗日纤维化的结构。

\section{特殊拉格朗日子流形的模空间}

真正理解 SYZ 猜想的第一步，是考察特殊拉格朗日纤维化的底上出现的结构。这些结构来自 McLean 关于特殊拉格朗日子流形模空间的定理 [45]，而这些结构及其相互关系由 Hitchin 在 [30] 中作了说明。我们在此概述其中一些想法。McLean 的定理说：紧 Calabi--Yau 流形 $X$ 中紧特殊拉格朗日子流形的形变模空间是无阻碍的，且 $M\subseteq X$（特殊拉格朗日）处的切空间典范同构于 $M$ 上调和 1-形式构成的空间。这一同构可以如下显式地看出。设 $\nu\in\Gamma(M,N_{M/X})$ 是 $M$ 在 $X$ 中的法向量场。那么 $(\iota(\nu)\omega)|_M$ 与 $(\iota(\nu)\operatorname{Im}\Omega)|_M$ 都被看成 $M$ 上良定义的微分形式：这需要把 $\nu$ 提升为向量场，但这一选择无关紧要，因为 $\omega$ 与 $\operatorname{Im}\Omega$ 在 $M$ 上限制为零。McLean 证明：若 $M$ 是特殊拉格朗日的，则
\[
\iota(\nu)\operatorname{Im}\Omega=-*\iota(\nu)\omega,
\]
其中 $*$ 表示 $M$ 上的 Hodge 星算子；并且 $\nu$ 对应于保持特殊拉格朗日条件的无穷小形变，当且仅当 $d(\iota(\nu)\omega)=d(\iota(\nu)\operatorname{Im}\Omega)=0$。这就给出了调和 1-形式与无穷小特殊拉格朗日形变之间的对应。

设 $f:X\to B$ 是纤维为环面的特殊拉格朗日纤维化，并暂时假设 $f$ 的所有纤维都是非奇异的。于是我们在 $B$ 上得到三个结构：两个仿射结构和一个度量，下面就会看到。

\begin{definition}{3.2.1}
设 $B$ 是 $n$ 维流形。$B$ 上的一个\emph{仿射结构}由一族坐标卡 $\{(U_i,\psi_i)\}$ 给出，其中 $\psi_i:U_i\to\R^n$，且其转移函数 $\psi_i\circ\psi_j^{-1}$ 属于 $\operatorname{Aff}(\R^n)$。若转移函数属于 $\R^n\rtimes GL(\Z^n)$，即其线性部分为整系数的，则称该仿射结构是\emph{热带的}（tropical）。若转移函数属于 $\operatorname{Aff}(\Z^n)$，则称该仿射结构是\emph{整的}（integral）。
\end{definition}

% [Tutor Note: 定义 2.1 中热带仿射结构与整仿射结构的区别在于平移：tropical 只要求线性部分整（转移函数在 $\R^n\rtimes GL(\Z^n)$ 中），integral 则要求平移也整（在 $\operatorname{Aff}(\Z^n)$ 中）。记号 $\rtimes$ 是半直积，不是普通的 $\times$；原文此处确为 $\rtimes$（pdftotext 文字层给出 U+22CA）。]

若仿射流形 $B$ 带有一个 Riemann 度量 $g$，且 $g$ 局部地由 $g_{ij}=\partial^2K/\partial y_i\partial y_j$ 给出，其中 $K$ 是某个凸函数、$y_1,\dots,y_n$ 是仿射坐标，则称该度量是\emph{仿射 Kähler 的}或 \emph{Hessian 的}。

于是我们如下得到这三个结构：

\para{仿射结构 1。}
对于 $f$ 的纤维 $X_b$ 的法向量场 $\nu$，$(\iota(\nu)\omega)|_{X_b}$ 是 $X_b$ 上良定义的 1-形式，其周期可以如下计算。设 $U\subseteq B$ 是小的开集，并设有子流形 $\gamma_1,\dots,\gamma_n\subseteq f^{-1}(U)$，它们是 $U$ 上的一族 1-循环，且使得对每个 $b\in U$，$\gamma_1\cap X_b,\dots,\gamma_n\cap X_b$ 构成 $H_1(X_b,\Z)$ 的一组基。考虑 $U$ 上的 1-形式 $\omega_1,\dots,\omega_n$，它们由
% [原文 p.7]
纤维上的积分定义：
\begin{equation*}
\omega_i(\nu)=\int_{X_b\cap\gamma_i}\iota(\nu)\omega,
\end{equation*}
其中 $\nu$ 是 $B$ 在 $b$ 处的切向量，我们可以把它提升为 $X_b$ 的法向量场。我们有 $\omega_i=f_*(\omega|_{\gamma_i})$，且由于 $\omega$ 是闭的，$\omega_i$ 也是闭的。于是存在 $U$ 上局部定义的函数 $y_1,\dots,y_n$，使得 $dy_i=\omega_i$。此外，这些函数在 $H_1(X_b,\Z)$ 的基的选取以及常数之下是良定义的。最后，由 McLean 定理，$\nu\mapsto\iota(\nu)\omega$ 给出 $T_{B,b}$ 与 $H^1(X_b,\R)$ 的同构，由此可见它们确实给出良定义的坐标。于是 $y_1,\dots,y_n$ 定义了 $B$ 上一个热带仿射结构的局部坐标。

% [Tutor Note: $\omega_i=f_*(\omega|_{\gamma_i})$ 中的 $f_*$ 是紧纤维支撑的推出（沿纤维积分）：先把 $\omega$ 限制到子流形 $\gamma_i$ 上，再对 $f$ 的纤维取积分，得到底 $B$ 上的 1-形式。这就是后文「纤维积分」的含义，不要与谱序列部分的 $R^pf_*$ 混淆。]

\para{仿射结构 2。}
用 $\operatorname{Im}\Omega$ 可以玩同样的把戏：选取子流形 $\Gamma_1,\dots,\Gamma_n\subseteq f^{-1}(U)$，它们是 $U$ 上 $n-1$ 维循环的族（即 $(n-1)$-循环族），且使得 $\Gamma_1\cap X_b,\dots,\Gamma_n\cap X_b$ 构成 $H^{n-1}(X_b,\Z)$ 的一组基。我们由 $\lambda_i=-f_*(\operatorname{Im}\Omega|_{\Gamma_i})$ 定义 $\lambda_i$，等价地，
\begin{equation*}
\lambda_i(\nu)=-\int_{X_b\cap\Gamma_i}\iota(\nu)\operatorname{Im}\Omega.
\end{equation*}
同样，$\lambda_1,\dots,\lambda_n$ 是闭 1-形式，局部地有 $\lambda_i=d\check y_i$，并且 $\check y_1,\dots,\check y_n$ 又是 $B$ 上一个热带仿射结构的仿射坐标。

% [Tutor Note: 两套坐标要区分：由 $\omega$ 得到 $dy_i=\omega_i$（仿射结构 1），由 $\operatorname{Im}\Omega$ 得到 $d\check y_i=\lambda_i$（仿射结构 2）；带 check 的 $\check y$ 对应 $\operatorname{Im}\Omega$ 那一侧，后文 Proposition 2.2 说明二者由 Legendre 变换相联系。]

\para{McLean 度量。}
$H^1(X_b,\R)$ 上的 Hodge 度量由
\begin{equation*}
g(\alpha,\beta)=\int_{X_b}\alpha\wedge*\beta
\end{equation*}
给出，其中 $\alpha,\beta$ 是调和 1-形式；它于是诱导出 $B$ 上的一个度量，该度量可以写成
\begin{equation*}
g(\nu_1,\nu_2)=-\int_{X_b}\iota(\nu_1)\omega\wedge\iota(\nu_2)\operatorname{Im}\Omega.
\end{equation*}

Hitchin [30] 的一个关键观察是：这些结构由 Legendre 变换相联系：

\begin{proposition}{3.2.2}
设 $y_1,\dots,y_n$ 是 $B$ 上关于由 $\omega$ 诱导的仿射结构的局部仿射坐标。则局部地存在 $B$ 上的函数 $K$，使得
\begin{equation*}
g(\partial/\partial y_i,\partial/\partial y_j)=\partial^2K/\partial y_i\partial y_j.
\end{equation*}
此外，$\check y_i=\partial K/\partial y_i$ 构成关于由 $\operatorname{Im}\Omega$ 诱导的仿射结构的一组仿射坐标；且若
\begin{equation*}
\check K(\check y_1,\dots,\check y_n)=\sum\check y_iy_i-K(y_1,\dots,y_n)
\end{equation*}
是 $K$ 的 Legendre 变换，则
\begin{equation*}
y_i=\partial\check K/\partial\check y_i
\end{equation*}
并且
\begin{equation*}
\partial^2\check K/\partial y_i\partial y_j=g(\partial/\partial\check y_i,\partial/\partial\check y_j).
\end{equation*}
\end{proposition}

% [Tutor Note: 命题末式原文把左边分母印作 $\partial y_i\partial y_j$（无 check），而右边度量取的是 $\partial/\partial\check y_i$、$\partial/\partial\check y_j$；按 Legendre 变换的对称性，左边分母应带 check，即 $\partial^2\check K/\partial\check y_i\partial\check y_j$。已用 1200 dpi 裁图（p.7，y≈596--618pt）复核，原文确实如此；正文照原文排，此为疑似原文笔误。]

\begin{proof}
取如上的族 $\gamma_1,\dots,\gamma_n,\Gamma_1,\dots,\Gamma_n$，定义在一个开邻域 $U$ 上，且两组基互为 Poincaré 对偶，即对 $b\in U$ 有 $(\gamma_i\cap X_b)\cdot(\Gamma_j\cap X_b)=\delta_{ij}$。设 $\gamma_1^*,\dots,\gamma_n^*$ 与 $\Gamma_1^*,\dots,\Gamma_n^*$ 分别是 $\Gamma(U,R^1f_*\Z)$ 与 $\Gamma(U,R^{n-1}f_*\Z)$ 的对偶基。由 $\gamma_i$ 的选取，我们得到局部坐标 $y_1,\dots,y_n$，使 $dy_i=\omega_i$，于是特别有
\begin{equation*}
\delta_{ij}=\omega_i(\partial/\partial y_j)=\int_{\gamma_i\cap X_b}\iota(\partial/\partial y_j)\omega,
\end{equation*}
